%% kapa: exam=1 round=37 year=2026 period=2 form=A subject=appraisal_law subject_ko=감정평가관계법규 seq=36 number=36
%% source: Q-Net 감정평가사 기출문제 (cst00309) artlSeq=5258924 file=raw/1008/5258924/2026년 제37회 감정평가사 1차 2교시.pdf
%% conversion: kordoc 4.24.51 -> tools/gen1.py (auto); review=proofread-sonnet; flags=none
\begin{question}{36}
도시 및 주거환경정비법령상 사업시행자가 사업시행계획인가를 받은 때에 의제될 수 있는 인ㆍ허가등에 해당하지 않는 것은?
\begin{choices}
  \item ｢토양환경보전법｣ 제12조에 따른 특정토양오염관리대상시설의 신고
  \item ｢하천법｣ 제33조에 따른 하천의 점용허가
  \item ｢유통산업발전법｣ 제8조에 따른 대규모점포등의 등록
  \item ｢도시공원 및 녹지 등에 관한 법률｣ 제16조의2에 따른 공원조성계획의 결정
  \item ｢전기안전관리법｣ 제8조에 따른 자가용전기설비의 공사계획의 인가 및 신고
\end{choices}
\end{question}
\answer{1}
